\section{Discussion}
\label{sec:limits}
Q7 verifies current-version protocol integration and lineage-checked
acquisition in two simulator scenes, but both exploratory datasets fail the
model gates and do not establish transfer. Historical Q4 used a shared
scripted scaffold and direct-node interface; profile checks used mock I/O.
Current profile and per-agent-binding adaptation, and the registered-harness
path combined with Q7, remain untested.

Q6 saved but did not apply a recipe; its cadence-bypassed stop does not prove
simultaneous sensor freshness. A separate default-paced check missed its
target and failed restoration, so neither trace establishes recipe-recovery
readiness. Approval records do not measure operator effectiveness
\cite{parasuramanmanzey2010complacency,parasuraman2000model,bainbridge1983ironies};
restricted tool access does not establish deployment security
\cite{greshake2023injection,liu2024pinjection,iec62443-3-3} or replace
independently engineered PLC/SIS protection \cite{iec61511}.

Further evaluation needs a versioned build, matched scenes, a predeclared
signal-pair and excitation plan, and default-paced recipe baselines. Physical
plant tests are required to establish transfer beyond repository simulators.

\newpage
\section{Conclusion}
\label{sec:conclusion}
AgentWorkShop integrates collaborative harnesses, production context, model
services, semantic I/O, and node/recipe workflows while retaining distinct
authorization, command-confirmation, and observation records. Historical
archives and the current two-scene run verify selected software paths, but
the current model gates withhold both recommendations. Successful
optimization, cross-scene transfer, cross-harness portability, recipe
recovery, plant transfer, and comparative control benefit remain unestablished.
